\documentclass[a4paper,10pt]{article}

% Package imports
\usepackage{latexsym}
\usepackage{xcolor}
\usepackage{float}
\usepackage{ragged2e}
\usepackage[empty]{fullpage}
\usepackage{wrapfig}
\usepackage{lipsum}
\usepackage{tabularx}
\usepackage{titlesec}
\usepackage{geometry}
\usepackage{marvosym}
\usepackage{verbatim}
\usepackage{enumitem}
\usepackage{fancyhdr}
\usepackage{multicol}
\usepackage{graphicx}
\usepackage{cfr-lm}
\usepackage[T1]{fontenc}
\usepackage{fontawesome5}

% Color definitions
\definecolor{darkblue}{RGB}{0,0,139}

% Page layout
\setlength{\multicolsep}{0pt}
\pagestyle{fancy}
\fancyhf{} % clear all header and footer fields
\fancyfoot{}
\renewcommand{\headrulewidth}{0pt}
\renewcommand{\footrulewidth}{0pt}
\geometry{left=1.4cm, top=0.8cm, right=1.2cm, bottom=1cm}
\setlength{\footskip}{5pt} % Addressing fancyhdr warning

% Hyperlink setup (moved after fancyhdr to address warning)
\usepackage[hidelinks]{hyperref}
\hypersetup{
    colorlinks=true,
    linkcolor=darkblue,
    filecolor=darkblue,
    urlcolor=darkblue,
}

% Custom box settings
\usepackage[most]{tcolorbox}
\tcbset{
    frame code={},
    center title,
    left=0pt,
    right=0pt,
    top=0pt,
    bottom=0pt,
    colback=gray!20,
    colframe=white,
    width=\dimexpr\textwidth\relax,
    enlarge left by=-2mm,
    boxsep=4pt,
    arc=0pt,outer arc=0pt,
}

% URL style
\urlstyle{same}

% Text alignment
\raggedright
\setlength{\tabcolsep}{0in}

% Section formatting
\titleformat{\section}{
  \vspace{-4pt}\scshape\raggedright\large
}{}{0em}{}[\color{black}\titlerule \vspace{-7pt}]

% Custom commands
\newcommand{\resumeItem}[2]{
  \item{
    \textbf{#1}{\hspace{0.5mm}#2 \vspace{-0.5mm}}
  }
}

\newcommand{\resumePOR}[3]{
\vspace{0.5mm}\item
    \begin{tabular*}{0.97\textwidth}[t]{l@{\extracolsep{\fill}}r}
        \textbf{#1}\hspace{0.3mm}#2 & \textit{\small{#3}}
    \end{tabular*}
    \vspace{-2mm}
}

\newcommand{\resumeSubheading}[4]{
\vspace{0.5mm}\item
    \begin{tabular*}{0.98\textwidth}[t]{l@{\extracolsep{\fill}}r}
        \textbf{#1} & \textit{\footnotesize{#4}} \\
        \textit{\footnotesize{#3}} &  \footnotesize{#2}\\
    \end{tabular*}
    \vspace{-2.4mm}
}

\newcommand{\resumeProject}[4]{
\vspace{0.5mm}\item
    \begin{tabular*}{0.98\textwidth}[t]{l@{\extracolsep{\fill}}r}
        \textbf{#1} & \textit{\footnotesize{#3}} \\
        \footnotesize{\textit{#2}} & \footnotesize{#4}
    \end{tabular*}
    \vspace{-2.4mm}
}

\newcommand{\resumeSubItem}[2]{\resumeItem{#1}{#2}\vspace{-4pt}}

\renewcommand{\labelitemi}{$\vcenter{\hbox{\tiny$\bullet$}}$}
\renewcommand{\labelitemii}{$\vcenter{\hbox{\tiny$\circ$}}$}

\newcommand{\resumeSubHeadingListStart}{\begin{itemize}[leftmargin=*,labelsep=1mm]}
\newcommand{\resumeHeadingSkillStart}{\begin{itemize}[leftmargin=*,itemsep=1.7mm, rightmargin=2ex]}
\newcommand{\resumeItemListStart}{\begin{itemize}[leftmargin=*,labelsep=1mm,itemsep=0.5mm]}

\newcommand{\resumeSubHeadingListEnd}{\end{itemize}\vspace{2mm}}
\newcommand{\resumeHeadingSkillEnd}{\end{itemize}\vspace{-2mm}}
\newcommand{\resumeItemListEnd}{\end{itemize}\vspace{-2mm}}
\newcommand{\cvsection}[1]{%
\vspace{2mm}
\begin{tcolorbox}
    \textbf{\large #1}
\end{tcolorbox}
    \vspace{-4mm}
}

\newcolumntype{L}{>{\raggedright\arraybackslash}X}%
\newcolumntype{R}{>{\raggedleft\arraybackslash}X}%
\newcolumntype{C}{>{\centering\arraybackslash}X}%

% Commands for icon sizing and positioning
\newcommand{\socialicon}[1]{\raisebox{-0.05em}{\resizebox{!}{1.2em}{#1}}}
\newcommand{\ieeeicon}[1]{\raisebox{-0.3em}{\resizebox{!}{1.3em}{#1}}}

% Font options
\newcommand{\headerfonti}{\fontfamily{phv}\selectfont} % Helvetica-like (similar to Arial/Calibri)
\newcommand{\headerfontii}{\fontfamily{ptm}\selectfont} % Times-like (similar to Times New Roman)
\newcommand{\headerfontiii}{\fontfamily{ppl}\selectfont} % Palatino (elegant serif)
\newcommand{\headerfontiv}{\fontfamily{pbk}\selectfont} % Bookman (readable serif)
\newcommand{\headerfontv}{\fontfamily{pag}\selectfont} % Avant Garde-like (similar to Trebuchet MS)
\newcommand{\headerfontvi}{\fontfamily{cmss}\selectfont} % Computer Modern Sans Serif
\newcommand{\headerfontvii}{\fontfamily{qhv}\selectfont} % Quasi-Helvetica (another Arial/Calibri alternative)
\newcommand{\headerfontviii}{\fontfamily{qpl}\selectfont} % Quasi-Palatino (another elegant serif option)
\newcommand{\headerfontix}{\fontfamily{qtm}\selectfont} % Quasi-Times (another Times New Roman alternative)
\newcommand{\headerfontx}{\fontfamily{bch}\selectfont} % Charter (clean serif font)

\begin{document}
\headerfontiii

% Header
\begin{center}
    {\Huge\textbf{SRINIVAS}}
\end{center}
\vspace{-4mm}

\begin{center}
    \small{
    +44-07553 994632 | \href{mailto:sxt589@student.bham.ac.uk}{sxt589@student.bham.ac.uk} |
    \href{https://srinivastb.netlify.app/}{Personal Website}
    }
\end{center}
\vspace{-6mm}

\begin{center}
    \small{
    \socialicon{\faLinkedin} \href{https://www.linkedin.com/in/srinivastb/}{/srinivastb} |
    \socialicon{\faGithub} \href{https://github.com/ch33nchan}{/ch33nchan} |
    \ieeeicon{\includegraphics[height=1.3em]{google-scholar-icon.png}} \href{https://scholar.google.com/citations?hl=en&view_op=list_works&gmla=AFix5Ma3lByqQf9mAXVpZaf8sSY5ZTWxMDpfNQg_F0y-3Nwvs4irFJjLqA_ey6Azdd0bfiNZb3hLC6i8Ibm3eA&user=UgR4sMUAAAAJ&inst=6297097477975492460}{/srinivas} |
    \socialicon{\faTwitter} \href{https://x.com/srinitwtts}{srinitwtts}
    }
\end{center}
\vspace{-4mm}
\begin{center}
    \small{Birmingham, United Kingdom}
\end{center}

\vspace{1mm}

\section{\textbf{About Me}}
\vspace{1mm}
\small{
\begin{itemize}
    \item Machine Learning Engineer with expertise in reinforcement learning and robotics. I blend research and engineering to build intelligent systems that solve real-world problems through a mix of theoretical innovation and practical implementation.
    
    \item I develop custom simulation environments and training frameworks that connect AI theory with robotics applications. My work in reward engineering and multi-objective optimization helps robots learn complex behaviors that transfer effectively to physical systems.
    
    \item My research combines reinforcement learning algorithms with software engineering best practices to create robust, scalable AI systems. I've optimized RL methods (PPO, SAC) that significantly reduce training time while improving performance in robotic control tasks.
    
    \item As a software engineer, I build end-to-end systems integrating perception, planning, and control components. I apply software architecture principles to ML systems, creating maintainable codebases that balance innovation with reliability for production environments.
\end{itemize}
}
\vspace{-2mm}
\section{\textbf{Technical Skills}}

\textbf{Research Areas}:  
Reinforcement Learning, Robotics Simulation, Multi-Agent Systems, Computer Vision, Natural Language Processing, Generative AI
\vspace{1mm}

\textbf{Robotics \& Simulation}:  
ROS2, Isaac Gym, Isaac Sim, MuJoCo, Gazebo, Unity ML-Agents, PyBullet, SLAM, Path Planning, Sensor Fusion
\vspace{1mm}

\textbf{ML Frameworks \& Algorithms}:  
PyTorch, TensorFlow, JAX, Stable Baselines3, PPO, SAC, TRPO, DDPG, TD3, A2C, DQN, Monte Carlo Tree Search
\vspace{1mm}

\textbf{Programming}:  
Python (Advanced), C++ (Intermediate), CUDA, SQL (Advanced), Bash, CMake
\vspace{1mm}

\textbf{Data Science \& Visualization}:  
NumPy, Pandas, Matplotlib, Seaborn, SciPy, Scikit-Learn, XGBoost, OpenCV, PCL, TensorBoard
\vspace{1mm}

\textbf{MLOps \& Infrastructure}:  
Docker, Kubernetes, Ray, Horovod, MLflow, Airflow, Git, CI/CD, TensorRT, ONNX
\vspace{1mm}

\textbf{Cloud \& Distributed Computing}:  
AWS (EC2, S3, SageMaker, BedRock), GCP (Vertex AI, Compute Engine), Multi-GPU Training, Distributed RL
\resumeHeadingSkillEnd


\section{\textbf{Education}}
\vspace{-0.4mm}
\resumeSubHeadingListStart

\resumeSubheading
{University of Birmingham}{Birmingham, United Kingdom}
{Masters in Artificial Intelligence \& Machine Learning}{Sept 2024 - Present}
\resumeItemListStart
\item \textbf{Global Masters Scholarship} - Awarded for academic excellence and leadership potential.
\item \textbf{Chancellor's Award} - Recognized for contributions to the university community.
\resumeItemListEnd

\resumeSubheading
{SRM Institute of Science \& Technology}{Chennai, India}
{Bachelors in Technology in Computer Science}{June 2024}
\resumeItemListStart
\item \textbf{First Class with Distinction} - Achieved a grade of 90.0\%.
\resumeItemListEnd

\resumeSubHeadingListEnd

\section{\textbf{Experience}}
\vspace{-0.4mm}
\resumeSubHeadingListStart

\resumeSubheading
{Alan Turing Institute [\href{https://research.birmingham.ac.uk/en/projects/human-machine-teaming}{\faIcon{globe}}]}{Birmingham, United Kingdom}
{\textbf{Research Associate} | Collaboration: University of Birmingham \& U.S. Army Research Institute}{Dec 2024 - Present}
  \item Developed new multi-agent reinforcement learning systems with PyTorch and PettingZoo for human-machine teaming. Used PPO to create information-asymmetric agents, improving coordination by 30\% in Stag Hunt scenarios through adaptive rewards.
  
  \item Created multi-agent systems with partial observability in PettingZoo, combining Monte Carlo Tree Search with Deep RL. Added attention mechanisms for agent communication, making policy convergence 40\% faster.
  
  \item Built a distributed training system for PyTorch models using Ray and Horovod, supporting 1000+ concurrent agents across multiple GPUs. Improved experience replay methods, making training twice as fast.
  
  \item Designed custom OpenAI Gym environments with adversarial training, using GANs for scenario generation. Made agents 45\% more resilient to unexpected situations through domain randomization.
\resumeItemListEnd

\clearpage  % 
\resumeSubheading
      {{RiskOpsAI [\href{https://www.optimeyes.ai/}{\faIcon{globe}}]}}{San Jose, Remote}
      {\textbf{AI ML Intern}}{Jun 2023 - Aug 2024}
      \resumeItemListStart
        \item Built a deep learning pipeline with TensorFlow and PyTorch, using ResNet-50 and Transformer models for classification. Optimized with TensorRT for 25\% faster inference and 15\% better accuracy.
        
        \item Created distributed ML infrastructure with Apache Airflow and MLflow, reducing training time by 20\% using Horovod on AWS GPU clusters. Set up automated data pipelines with feature engineering.
        
        \item Developed a predictive analytics system with PostgreSQL and BigQuery, using scikit-learn and XGBoost for imbalanced datasets. Improved decision-making efficiency by 30\% through query optimization.
      \resumeItemListEnd
    \resumeSubheading
      {{SRM Institute of Science \& Technology [\href{https://www.srmist.edu.in/faculty/vaishnavi-moorthy/}{\faIcon{globe}}]}}{Chennai, India}
      {\textbf{Researcher under Dr. Vaishnavi Moorthy}}{Oct 2022 - Feb 2024}
      \resumeItemListStart
        \item Developed an autonomous navigation system using ROS2, combining LiDAR, RGB-D, and IMU sensors with Extended Kalman Filters. Improved localization by 15\% through sensor fusion adjustments.
        
        \item Created a SLAM system with SAC and TRPO algorithms in PyTorch, improving path planning with RRT* and A*. Reduced navigation errors by 25\% with dynamic obstacle avoidance.
        
        \item Built a real-time perception pipeline using OpenCV and PCL, adding YOLOv7 for object detection. Achieved 20ms latency and 95\% detection accuracy in changing environments.
      \resumeItemListEnd

\section{\textbf{Projects}}
 \resumeSubheading
   {RLlama: Memory-Augmented RL Framework for LLMs}
   {\href{https://github.com/ch33nchan/RLlama}{\textcolor{darkblue}{\faGithub}}
    \href{https://pypi.org/project/rllama/}{\textcolor{blue}{\faPython}}
    \href{https://news.ycombinator.com/item?id=43167073}{\textcolor{orange}{\faHackerNews}}
   }
   {Tools: Python, PyTorch, Transformers, Accelerate, PEFT, Optuna, Reinforcement Learning}
   {Jan 2025 - Present}
   \resumeItemListStart
     \item Built a modular reward shaping framework with composable components, dynamic weight scheduling, and a component registry, giving fine control over agent learning signals.
     
     \item Created an automated reward parameter optimization system using Bayesian methods with Optuna and visualization tools to speed up hyperparameter tuning.
     
     \item Developed an episodic memory system with hierarchical importance sampling, achieving 87\% memory compression and reducing LLM context needs by 65\% while maintaining 93\% task performance.
     
     \item Integrated multiple RL algorithms in the agent core, making convergence 42\% faster and improving sample efficiency by 35\% compared to standard implementations.
     
     \item Built a distributed training pipeline with Accelerate, scaling to 32-node clusters and processing over 1M tokens/sec, reducing training time by 78\%.
 \resumeItemListEnd

  \resumeSubheading
    {Lightning Quant: Modular Trading Strategy Framework}{\href{https://github.com/ch33nchan/light-quant-main}{\textcolor{darkblue}{\faGithub}}}
    {Tools: Python, PyTorch Lightning, Pandas, NumPy, Scikit-learn, TA-Lib}{Jul 2023 - Oct 2023}
    \resumeItemListStart
      \item Created a modular quantitative trading pipeline using PyTorch Lightning, covering data acquisition, feature engineering, strategy optimization, and model training.
      
      \item Built a feature engineering module generating 20+ technical indicators (NATR, Aroon, RSI, VWAP) with expanding windows and ranking, stored efficiently using Parquet format.
      
      \item Developed a brute-force optimization engine for strategy parameters (like dual moving averages), evaluating performance based on CAGR, Sharpe Ratio, and Max Drawdown.
  \resumeItemListEnd

  \resumeSubheading
    {TinyGrad RLCV: Lightweight RL for Computer Vision}{\href{https://github.com/ch33nchan/tinygrad-rlcv}{\textcolor{darkblue}{\faGithub}}}
    {Tools: Python, TinyGrad, OpenCV, NumPy, Reinforcement Learning}{Oct 2023 - Dec 2023}
    \resumeItemListStart
      \item Created ultra-lightweight neural network operations with TinyGrad, achieving a small memory footprint (<10MB) for deployment on resource-limited edge devices (Mobile CPUs, ARM).
      
      \item Built a memory-optimized DQN agent with prioritized experience replay and flexible policy networks, enabling stable learning and real-time inference (>30 FPS) on target hardware.
      
      \item Developed an optimized OpenCV pipeline for efficient, real-time feature extraction from webcam streams, enabling robust object tracking on low-power systems.
    \resumeItemListEnd

  \resumeSubHeadingListEnd

\vspace{-2mm}
\section{\textbf{Publications} \hfill \textcolor{darkblue}{\scriptsize C=Conference, J=Journal, P=Patent, S=In Submission, T=Thesis}}
\vspace{0.1mm}
\small{
\begin{enumerate}[leftmargin=*, labelsep=0.5em, align=left, widest={[\textbf{S.1}]}, itemindent=0em, label={\textbf{[\arabic*]}]}]
\item[\textbf{[C.1]}] V. Moorthy, T. B. Srinivas. (2024). \href{https://doi.org/10.1007/978-981-97-1329-5_36}{\textbf{Enhancing Low-Light Surveillance Images with MirNet: A Keras-Based Approach}}. In \textit{Smart Trends in Computing and Communications}, pp. 445–458. Springer. 2024, Location SmartCom Pune, India. DOI: \href{https://doi.org/10.1007/978-981-97-1329-5_36}{10.1007/978-981-97-1329-5_36}.
\item[\textbf{[C.2]}] V. Moorthy, D. Shah, S. Kapoor, S. T. B., A. Lal. (2023). \href{https://doi.org/10.1109/CONECCT58849.2023.10234740}{\textbf{ALATS: Analysis of Localization Algorithms in Terrestrial Surveillance Bots}}. In \textit{2023 IEEE International Conference on Electronics, Computing and Communication Technologies (CONECCT)}, pp. 1–6. IEEE. 2023, Bangalore, India. DOI: \href{https://doi.org/10.1109/CONECCT58849.2023.10234740}{10.1109/CONECCT58849.2023.10234740}.
\end{enumerate}
}

\end{document}
